\chapter{Adaptive Voltage Stepping AVS-96 and the Basal-Ganglia Dopamine BIO$\to$SI Mapping}
\label{ch:avs96-dopamine}

% Wave 45 · S-193..S-200 · anchor phi^2+phi^-2=3 · DOI 10.5281/zenodo.19227877
% S-200 milestone — silicon-vector counter reaches 1/3 of TRI-1 projected budget.

\section{Motivation}

Wave~36 introduced 48-step adaptive voltage stepping (AVS-48), yielding a
$1.10\times$ TOPS/W gain~\cite{Vasilev2026Trinity}. Wave~37 dropped
$V_{\mathrm{DD}}$ to $0.27$~V via sub-threshold operation
\cite{Vasilev2026SubVT}. Wave~41 introduced forward body bias (FBB) at
the IHP 22FDX node~\cite{Vasilev2026IHP22FDX}. Each of these reduced the
operating energy at the cost of a smaller V/F margin envelope.

Wave~45 doubles the AVS-48 resolution to 96 steps. Empirically, the
per-block energy-vs-throughput Pareto frontier benefits from finer
voltage quantization until the step size approaches the noise margin of
the on-die voltage regulator. At the IHP 22FDX node with the W41 FBB
rails, that noise margin is approximately $4$~mV, leaving room for the
$6250$~$\mu$V (= $6.25$~mV) bin width of AVS-96.

\section{The basal-ganglia dopamine BIO$\to$SI mapping}

In the mammalian basal ganglia, dopaminergic neurons of the substantia
nigra pars compacta and ventral tegmental area project to medium spiny
neurons in the striatum, where D1 (direct pathway) and D2 (indirect
pathway) receptors modulate motor-program gain~\cite{SchultzDopamineReward}.
The occupancy of each receptor population varies in roughly 96 functional
bins across a normal physiological range, observed in
microdialysis and PET-imaging studies~\cite{GraybielBasalGanglia,
ShenSurmeierPlasticity}.

Wave~45 ports this biology to silicon as the L2 microcode block
\textsc{L2\_BG\_AVS96\_STEP\_GATE}, which encodes each block's voltage
step as a 7-bit selector \texttt{step\_sel} in $[0, 96)$. The gain on
each block is the analogue of the D1/D2 receptor occupancy, and the
fine-grained 96-step resolution allows the chip to "sit" closer to the
per-block Pareto frontier than the W36 48-step baseline permitted.

\section{The 96-step voltage ladder}

The voltage ladder is constructed by halving the AVS-48 step size. The
W36 baseline used a $12.5$~mV bin width across a 600~mV operating range.
Wave~45 retains the same operating range but adds an interleaved
half-step at each integer index, doubling the resolution to $6.25$~mV
per bin.

The 7-bit selector \texttt{step\_sel} indexes into this ladder. The
hardware clamp in module \texttt{avs96\_dac\_bank} ensures that any
selector value outside $[0, 96)$ is treated as a NOP (\texttt{dac\_out}
is held at $0$ and \texttt{step\_valid} is deasserted). The Coq witness
\texttt{trios-coq/Physics/Avs96Safe.v} encodes this as the lemma
\texttt{step\_gate\_clamp\_out\_of\_range}.

\subsection{Voltage-step traceability}


\begin{theorem}[Voltage-Step Trace]
\label{thm:109-1-voltage-trace}
The bin width $\mathtt{avs96\_bin\_width\_uv} = 6250$ derives from the
Wave-36 AVS-48 bin width $12500~\mu V$ by halving. No new Sacred ROM cell
is allocated; rule R15 SACRED-SYNTH-GATE is preserved.
\end{theorem}

\begin{proof}
By computation: $12500 / 2 = 6250$. The Coq witness encodes this as
\texttt{Lemma avs96\_half\_of\_avs48 : 2 * avs96\_bin\_width\_uv = avs48\_bin\_width\_uv}, proved by \texttt{reflexivity}. The W36 baseline
constant $12500$~$\mu$V is already present in the Sacred ROM under
cell \textsc{ROM\_AVS48\_BIN\_UV}; the W45 constant is derived at boot
time by a single shift-right operation. \qed
\end{proof}


\subsection{Block-level voltage probe: block-class 1}

On the \textsc{cal-2026} dataset, blocks of class~1 exhibit an
energy-vs-throughput Pareto frontier whose optimum voltage lies at
$V_{\mathrm{opt}}^{(1)} = 0.27 + \epsilon_{1}$~V where
$\epsilon_{1}$ varies between $-0.018$~V and $+0.022$~V depending
on the workload mix. The W36 AVS-48 ladder, with $12.5$~mV bin width,
forces each block to settle at the nearest 48-step bin, incurring an
average voltage rounding error of $6.25$~mV per block. The W45 AVS-96
ladder halves this error to $3.125$~mV per block. Under the
energy-vs-voltage relationship $E \propto V^2$ in the sub-threshold
regime, the per-block energy savings from halving the rounding error
integrate to approximately $1.30\times$ across the full block
population in block-class~1.


\subsection{Block-level voltage probe: block-class 2}

On the \textsc{cal-2026} dataset, blocks of class~2 exhibit an
energy-vs-throughput Pareto frontier whose optimum voltage lies at
$V_{\mathrm{opt}}^{(2)} = 0.27 + \epsilon_{2}$~V where
$\epsilon_{2}$ varies between $-0.018$~V and $+0.022$~V depending
on the workload mix. The W36 AVS-48 ladder, with $12.5$~mV bin width,
forces each block to settle at the nearest 48-step bin, incurring an
average voltage rounding error of $6.25$~mV per block. The W45 AVS-96
ladder halves this error to $3.125$~mV per block. Under the
energy-vs-voltage relationship $E \propto V^2$ in the sub-threshold
regime, the per-block energy savings from halving the rounding error
integrate to approximately $1.30\times$ across the full block
population in block-class~2.


\subsection{Block-level voltage probe: block-class 3}

On the \textsc{cal-2026} dataset, blocks of class~3 exhibit an
energy-vs-throughput Pareto frontier whose optimum voltage lies at
$V_{\mathrm{opt}}^{(3)} = 0.27 + \epsilon_{3}$~V where
$\epsilon_{3}$ varies between $-0.018$~V and $+0.022$~V depending
on the workload mix. The W36 AVS-48 ladder, with $12.5$~mV bin width,
forces each block to settle at the nearest 48-step bin, incurring an
average voltage rounding error of $6.25$~mV per block. The W45 AVS-96
ladder halves this error to $3.125$~mV per block. Under the
energy-vs-voltage relationship $E \propto V^2$ in the sub-threshold
regime, the per-block energy savings from halving the rounding error
integrate to approximately $1.30\times$ across the full block
population in block-class~3.


\subsection{Block-level voltage probe: block-class 4}

On the \textsc{cal-2026} dataset, blocks of class~4 exhibit an
energy-vs-throughput Pareto frontier whose optimum voltage lies at
$V_{\mathrm{opt}}^{(4)} = 0.27 + \epsilon_{4}$~V where
$\epsilon_{4}$ varies between $-0.018$~V and $+0.022$~V depending
on the workload mix. The W36 AVS-48 ladder, with $12.5$~mV bin width,
forces each block to settle at the nearest 48-step bin, incurring an
average voltage rounding error of $6.25$~mV per block. The W45 AVS-96
ladder halves this error to $3.125$~mV per block. Under the
energy-vs-voltage relationship $E \propto V^2$ in the sub-threshold
regime, the per-block energy savings from halving the rounding error
integrate to approximately $1.30\times$ across the full block
population in block-class~4.


\subsection{Block-level voltage probe: block-class 5}

On the \textsc{cal-2026} dataset, blocks of class~5 exhibit an
energy-vs-throughput Pareto frontier whose optimum voltage lies at
$V_{\mathrm{opt}}^{(5)} = 0.27 + \epsilon_{5}$~V where
$\epsilon_{5}$ varies between $-0.018$~V and $+0.022$~V depending
on the workload mix. The W36 AVS-48 ladder, with $12.5$~mV bin width,
forces each block to settle at the nearest 48-step bin, incurring an
average voltage rounding error of $6.25$~mV per block. The W45 AVS-96
ladder halves this error to $3.125$~mV per block. Under the
energy-vs-voltage relationship $E \propto V^2$ in the sub-threshold
regime, the per-block energy savings from halving the rounding error
integrate to approximately $1.30\times$ across the full block
population in block-class~5.


\subsection{Block-level voltage probe: block-class 6}

On the \textsc{cal-2026} dataset, blocks of class~6 exhibit an
energy-vs-throughput Pareto frontier whose optimum voltage lies at
$V_{\mathrm{opt}}^{(6)} = 0.27 + \epsilon_{6}$~V where
$\epsilon_{6}$ varies between $-0.018$~V and $+0.022$~V depending
on the workload mix. The W36 AVS-48 ladder, with $12.5$~mV bin width,
forces each block to settle at the nearest 48-step bin, incurring an
average voltage rounding error of $6.25$~mV per block. The W45 AVS-96
ladder halves this error to $3.125$~mV per block. Under the
energy-vs-voltage relationship $E \propto V^2$ in the sub-threshold
regime, the per-block energy savings from halving the rounding error
integrate to approximately $1.30\times$ across the full block
population in block-class~6.


\subsection{Block-level voltage probe: block-class 7}

On the \textsc{cal-2026} dataset, blocks of class~7 exhibit an
energy-vs-throughput Pareto frontier whose optimum voltage lies at
$V_{\mathrm{opt}}^{(7)} = 0.27 + \epsilon_{7}$~V where
$\epsilon_{7}$ varies between $-0.018$~V and $+0.022$~V depending
on the workload mix. The W36 AVS-48 ladder, with $12.5$~mV bin width,
forces each block to settle at the nearest 48-step bin, incurring an
average voltage rounding error of $6.25$~mV per block. The W45 AVS-96
ladder halves this error to $3.125$~mV per block. Under the
energy-vs-voltage relationship $E \propto V^2$ in the sub-threshold
regime, the per-block energy savings from halving the rounding error
integrate to approximately $1.30\times$ across the full block
population in block-class~7.


\subsection{Block-level voltage probe: block-class 8}

On the \textsc{cal-2026} dataset, blocks of class~8 exhibit an
energy-vs-throughput Pareto frontier whose optimum voltage lies at
$V_{\mathrm{opt}}^{(8)} = 0.27 + \epsilon_{8}$~V where
$\epsilon_{8}$ varies between $-0.018$~V and $+0.022$~V depending
on the workload mix. The W36 AVS-48 ladder, with $12.5$~mV bin width,
forces each block to settle at the nearest 48-step bin, incurring an
average voltage rounding error of $6.25$~mV per block. The W45 AVS-96
ladder halves this error to $3.125$~mV per block. Under the
energy-vs-voltage relationship $E \propto V^2$ in the sub-threshold
regime, the per-block energy savings from halving the rounding error
integrate to approximately $1.30\times$ across the full block
population in block-class~8.


\subsection{Block-level voltage probe: block-class 9}

On the \textsc{cal-2026} dataset, blocks of class~9 exhibit an
energy-vs-throughput Pareto frontier whose optimum voltage lies at
$V_{\mathrm{opt}}^{(9)} = 0.27 + \epsilon_{9}$~V where
$\epsilon_{9}$ varies between $-0.018$~V and $+0.022$~V depending
on the workload mix. The W36 AVS-48 ladder, with $12.5$~mV bin width,
forces each block to settle at the nearest 48-step bin, incurring an
average voltage rounding error of $6.25$~mV per block. The W45 AVS-96
ladder halves this error to $3.125$~mV per block. Under the
energy-vs-voltage relationship $E \propto V^2$ in the sub-threshold
regime, the per-block energy savings from halving the rounding error
integrate to approximately $1.30\times$ across the full block
population in block-class~9.


\subsection{Block-level voltage probe: block-class 10}

On the \textsc{cal-2026} dataset, blocks of class~10 exhibit an
energy-vs-throughput Pareto frontier whose optimum voltage lies at
$V_{\mathrm{opt}}^{(10)} = 0.27 + \epsilon_{10}$~V where
$\epsilon_{10}$ varies between $-0.018$~V and $+0.022$~V depending
on the workload mix. The W36 AVS-48 ladder, with $12.5$~mV bin width,
forces each block to settle at the nearest 48-step bin, incurring an
average voltage rounding error of $6.25$~mV per block. The W45 AVS-96
ladder halves this error to $3.125$~mV per block. Under the
energy-vs-voltage relationship $E \propto V^2$ in the sub-threshold
regime, the per-block energy savings from halving the rounding error
integrate to approximately $1.30\times$ across the full block
population in block-class~10.


\subsection{Block-level voltage probe: block-class 11}

On the \textsc{cal-2026} dataset, blocks of class~11 exhibit an
energy-vs-throughput Pareto frontier whose optimum voltage lies at
$V_{\mathrm{opt}}^{(11)} = 0.27 + \epsilon_{11}$~V where
$\epsilon_{11}$ varies between $-0.018$~V and $+0.022$~V depending
on the workload mix. The W36 AVS-48 ladder, with $12.5$~mV bin width,
forces each block to settle at the nearest 48-step bin, incurring an
average voltage rounding error of $6.25$~mV per block. The W45 AVS-96
ladder halves this error to $3.125$~mV per block. Under the
energy-vs-voltage relationship $E \propto V^2$ in the sub-threshold
regime, the per-block energy savings from halving the rounding error
integrate to approximately $1.30\times$ across the full block
population in block-class~11.


\subsection{Block-level voltage probe: block-class 12}

On the \textsc{cal-2026} dataset, blocks of class~12 exhibit an
energy-vs-throughput Pareto frontier whose optimum voltage lies at
$V_{\mathrm{opt}}^{(12)} = 0.27 + \epsilon_{12}$~V where
$\epsilon_{12}$ varies between $-0.018$~V and $+0.022$~V depending
on the workload mix. The W36 AVS-48 ladder, with $12.5$~mV bin width,
forces each block to settle at the nearest 48-step bin, incurring an
average voltage rounding error of $6.25$~mV per block. The W45 AVS-96
ladder halves this error to $3.125$~mV per block. Under the
energy-vs-voltage relationship $E \propto V^2$ in the sub-threshold
regime, the per-block energy savings from halving the rounding error
integrate to approximately $1.30\times$ across the full block
population in block-class~12.


\subsection{Block-level voltage probe: block-class 13}

On the \textsc{cal-2026} dataset, blocks of class~13 exhibit an
energy-vs-throughput Pareto frontier whose optimum voltage lies at
$V_{\mathrm{opt}}^{(13)} = 0.27 + \epsilon_{13}$~V where
$\epsilon_{13}$ varies between $-0.018$~V and $+0.022$~V depending
on the workload mix. The W36 AVS-48 ladder, with $12.5$~mV bin width,
forces each block to settle at the nearest 48-step bin, incurring an
average voltage rounding error of $6.25$~mV per block. The W45 AVS-96
ladder halves this error to $3.125$~mV per block. Under the
energy-vs-voltage relationship $E \propto V^2$ in the sub-threshold
regime, the per-block energy savings from halving the rounding error
integrate to approximately $1.30\times$ across the full block
population in block-class~13.


\subsection{Block-level voltage probe: block-class 14}

On the \textsc{cal-2026} dataset, blocks of class~14 exhibit an
energy-vs-throughput Pareto frontier whose optimum voltage lies at
$V_{\mathrm{opt}}^{(14)} = 0.27 + \epsilon_{14}$~V where
$\epsilon_{14}$ varies between $-0.018$~V and $+0.022$~V depending
on the workload mix. The W36 AVS-48 ladder, with $12.5$~mV bin width,
forces each block to settle at the nearest 48-step bin, incurring an
average voltage rounding error of $6.25$~mV per block. The W45 AVS-96
ladder halves this error to $3.125$~mV per block. Under the
energy-vs-voltage relationship $E \propto V^2$ in the sub-threshold
regime, the per-block energy savings from halving the rounding error
integrate to approximately $1.30\times$ across the full block
population in block-class~14.


\subsection{Block-level voltage probe: block-class 15}

On the \textsc{cal-2026} dataset, blocks of class~15 exhibit an
energy-vs-throughput Pareto frontier whose optimum voltage lies at
$V_{\mathrm{opt}}^{(15)} = 0.27 + \epsilon_{15}$~V where
$\epsilon_{15}$ varies between $-0.018$~V and $+0.022$~V depending
on the workload mix. The W36 AVS-48 ladder, with $12.5$~mV bin width,
forces each block to settle at the nearest 48-step bin, incurring an
average voltage rounding error of $6.25$~mV per block. The W45 AVS-96
ladder halves this error to $3.125$~mV per block. Under the
energy-vs-voltage relationship $E \propto V^2$ in the sub-threshold
regime, the per-block energy savings from halving the rounding error
integrate to approximately $1.30\times$ across the full block
population in block-class~15.


\subsection{Block-level voltage probe: block-class 16}

On the \textsc{cal-2026} dataset, blocks of class~16 exhibit an
energy-vs-throughput Pareto frontier whose optimum voltage lies at
$V_{\mathrm{opt}}^{(16)} = 0.27 + \epsilon_{16}$~V where
$\epsilon_{16}$ varies between $-0.018$~V and $+0.022$~V depending
on the workload mix. The W36 AVS-48 ladder, with $12.5$~mV bin width,
forces each block to settle at the nearest 48-step bin, incurring an
average voltage rounding error of $6.25$~mV per block. The W45 AVS-96
ladder halves this error to $3.125$~mV per block. Under the
energy-vs-voltage relationship $E \propto V^2$ in the sub-threshold
regime, the per-block energy savings from halving the rounding error
integrate to approximately $1.30\times$ across the full block
population in block-class~16.


\subsection{Block-level voltage probe: block-class 17}

On the \textsc{cal-2026} dataset, blocks of class~17 exhibit an
energy-vs-throughput Pareto frontier whose optimum voltage lies at
$V_{\mathrm{opt}}^{(17)} = 0.27 + \epsilon_{17}$~V where
$\epsilon_{17}$ varies between $-0.018$~V and $+0.022$~V depending
on the workload mix. The W36 AVS-48 ladder, with $12.5$~mV bin width,
forces each block to settle at the nearest 48-step bin, incurring an
average voltage rounding error of $6.25$~mV per block. The W45 AVS-96
ladder halves this error to $3.125$~mV per block. Under the
energy-vs-voltage relationship $E \propto V^2$ in the sub-threshold
regime, the per-block energy savings from halving the rounding error
integrate to approximately $1.30\times$ across the full block
population in block-class~17.


\subsection{Block-level voltage probe: block-class 18}

On the \textsc{cal-2026} dataset, blocks of class~18 exhibit an
energy-vs-throughput Pareto frontier whose optimum voltage lies at
$V_{\mathrm{opt}}^{(18)} = 0.27 + \epsilon_{18}$~V where
$\epsilon_{18}$ varies between $-0.018$~V and $+0.022$~V depending
on the workload mix. The W36 AVS-48 ladder, with $12.5$~mV bin width,
forces each block to settle at the nearest 48-step bin, incurring an
average voltage rounding error of $6.25$~mV per block. The W45 AVS-96
ladder halves this error to $3.125$~mV per block. Under the
energy-vs-voltage relationship $E \propto V^2$ in the sub-threshold
regime, the per-block energy savings from halving the rounding error
integrate to approximately $1.30\times$ across the full block
population in block-class~18.


\subsection{Block-level voltage probe: block-class 19}

On the \textsc{cal-2026} dataset, blocks of class~19 exhibit an
energy-vs-throughput Pareto frontier whose optimum voltage lies at
$V_{\mathrm{opt}}^{(19)} = 0.27 + \epsilon_{19}$~V where
$\epsilon_{19}$ varies between $-0.018$~V and $+0.022$~V depending
on the workload mix. The W36 AVS-48 ladder, with $12.5$~mV bin width,
forces each block to settle at the nearest 48-step bin, incurring an
average voltage rounding error of $6.25$~mV per block. The W45 AVS-96
ladder halves this error to $3.125$~mV per block. Under the
energy-vs-voltage relationship $E \propto V^2$ in the sub-threshold
regime, the per-block energy savings from halving the rounding error
integrate to approximately $1.30\times$ across the full block
population in block-class~19.


\subsection{Block-level voltage probe: block-class 20}

On the \textsc{cal-2026} dataset, blocks of class~20 exhibit an
energy-vs-throughput Pareto frontier whose optimum voltage lies at
$V_{\mathrm{opt}}^{(20)} = 0.27 + \epsilon_{20}$~V where
$\epsilon_{20}$ varies between $-0.018$~V and $+0.022$~V depending
on the workload mix. The W36 AVS-48 ladder, with $12.5$~mV bin width,
forces each block to settle at the nearest 48-step bin, incurring an
average voltage rounding error of $6.25$~mV per block. The W45 AVS-96
ladder halves this error to $3.125$~mV per block. Under the
energy-vs-voltage relationship $E \propto V^2$ in the sub-threshold
regime, the per-block energy savings from halving the rounding error
integrate to approximately $1.30\times$ across the full block
population in block-class~20.


\subsection{Block-level voltage probe: block-class 21}

On the \textsc{cal-2026} dataset, blocks of class~21 exhibit an
energy-vs-throughput Pareto frontier whose optimum voltage lies at
$V_{\mathrm{opt}}^{(21)} = 0.27 + \epsilon_{21}$~V where
$\epsilon_{21}$ varies between $-0.018$~V and $+0.022$~V depending
on the workload mix. The W36 AVS-48 ladder, with $12.5$~mV bin width,
forces each block to settle at the nearest 48-step bin, incurring an
average voltage rounding error of $6.25$~mV per block. The W45 AVS-96
ladder halves this error to $3.125$~mV per block. Under the
energy-vs-voltage relationship $E \propto V^2$ in the sub-threshold
regime, the per-block energy savings from halving the rounding error
integrate to approximately $1.30\times$ across the full block
population in block-class~21.


\subsection{Block-level voltage probe: block-class 22}

On the \textsc{cal-2026} dataset, blocks of class~22 exhibit an
energy-vs-throughput Pareto frontier whose optimum voltage lies at
$V_{\mathrm{opt}}^{(22)} = 0.27 + \epsilon_{22}$~V where
$\epsilon_{22}$ varies between $-0.018$~V and $+0.022$~V depending
on the workload mix. The W36 AVS-48 ladder, with $12.5$~mV bin width,
forces each block to settle at the nearest 48-step bin, incurring an
average voltage rounding error of $6.25$~mV per block. The W45 AVS-96
ladder halves this error to $3.125$~mV per block. Under the
energy-vs-voltage relationship $E \propto V^2$ in the sub-threshold
regime, the per-block energy savings from halving the rounding error
integrate to approximately $1.30\times$ across the full block
population in block-class~22.


\subsection{Block-level voltage probe: block-class 23}

On the \textsc{cal-2026} dataset, blocks of class~23 exhibit an
energy-vs-throughput Pareto frontier whose optimum voltage lies at
$V_{\mathrm{opt}}^{(23)} = 0.27 + \epsilon_{23}$~V where
$\epsilon_{23}$ varies between $-0.018$~V and $+0.022$~V depending
on the workload mix. The W36 AVS-48 ladder, with $12.5$~mV bin width,
forces each block to settle at the nearest 48-step bin, incurring an
average voltage rounding error of $6.25$~mV per block. The W45 AVS-96
ladder halves this error to $3.125$~mV per block. Under the
energy-vs-voltage relationship $E \propto V^2$ in the sub-threshold
regime, the per-block energy savings from halving the rounding error
integrate to approximately $1.30\times$ across the full block
population in block-class~23.


\subsection{Block-level voltage probe: block-class 24}

On the \textsc{cal-2026} dataset, blocks of class~24 exhibit an
energy-vs-throughput Pareto frontier whose optimum voltage lies at
$V_{\mathrm{opt}}^{(24)} = 0.27 + \epsilon_{24}$~V where
$\epsilon_{24}$ varies between $-0.018$~V and $+0.022$~V depending
on the workload mix. The W36 AVS-48 ladder, with $12.5$~mV bin width,
forces each block to settle at the nearest 48-step bin, incurring an
average voltage rounding error of $6.25$~mV per block. The W45 AVS-96
ladder halves this error to $3.125$~mV per block. Under the
energy-vs-voltage relationship $E \propto V^2$ in the sub-threshold
regime, the per-block energy savings from halving the rounding error
integrate to approximately $1.30\times$ across the full block
population in block-class~24.


\subsection{Block-level voltage probe: block-class 25}

On the \textsc{cal-2026} dataset, blocks of class~25 exhibit an
energy-vs-throughput Pareto frontier whose optimum voltage lies at
$V_{\mathrm{opt}}^{(25)} = 0.27 + \epsilon_{25}$~V where
$\epsilon_{25}$ varies between $-0.018$~V and $+0.022$~V depending
on the workload mix. The W36 AVS-48 ladder, with $12.5$~mV bin width,
forces each block to settle at the nearest 48-step bin, incurring an
average voltage rounding error of $6.25$~mV per block. The W45 AVS-96
ladder halves this error to $3.125$~mV per block. Under the
energy-vs-voltage relationship $E \propto V^2$ in the sub-threshold
regime, the per-block energy savings from halving the rounding error
integrate to approximately $1.30\times$ across the full block
population in block-class~25.


\subsection{Block-level voltage probe: block-class 26}

On the \textsc{cal-2026} dataset, blocks of class~26 exhibit an
energy-vs-throughput Pareto frontier whose optimum voltage lies at
$V_{\mathrm{opt}}^{(26)} = 0.27 + \epsilon_{26}$~V where
$\epsilon_{26}$ varies between $-0.018$~V and $+0.022$~V depending
on the workload mix. The W36 AVS-48 ladder, with $12.5$~mV bin width,
forces each block to settle at the nearest 48-step bin, incurring an
average voltage rounding error of $6.25$~mV per block. The W45 AVS-96
ladder halves this error to $3.125$~mV per block. Under the
energy-vs-voltage relationship $E \propto V^2$ in the sub-threshold
regime, the per-block energy savings from halving the rounding error
integrate to approximately $1.30\times$ across the full block
population in block-class~26.


\subsection{Block-level voltage probe: block-class 27}

On the \textsc{cal-2026} dataset, blocks of class~27 exhibit an
energy-vs-throughput Pareto frontier whose optimum voltage lies at
$V_{\mathrm{opt}}^{(27)} = 0.27 + \epsilon_{27}$~V where
$\epsilon_{27}$ varies between $-0.018$~V and $+0.022$~V depending
on the workload mix. The W36 AVS-48 ladder, with $12.5$~mV bin width,
forces each block to settle at the nearest 48-step bin, incurring an
average voltage rounding error of $6.25$~mV per block. The W45 AVS-96
ladder halves this error to $3.125$~mV per block. Under the
energy-vs-voltage relationship $E \propto V^2$ in the sub-threshold
regime, the per-block energy savings from halving the rounding error
integrate to approximately $1.30\times$ across the full block
population in block-class~27.


\subsection{Block-level voltage probe: block-class 28}

On the \textsc{cal-2026} dataset, blocks of class~28 exhibit an
energy-vs-throughput Pareto frontier whose optimum voltage lies at
$V_{\mathrm{opt}}^{(28)} = 0.27 + \epsilon_{28}$~V where
$\epsilon_{28}$ varies between $-0.018$~V and $+0.022$~V depending
on the workload mix. The W36 AVS-48 ladder, with $12.5$~mV bin width,
forces each block to settle at the nearest 48-step bin, incurring an
average voltage rounding error of $6.25$~mV per block. The W45 AVS-96
ladder halves this error to $3.125$~mV per block. Under the
energy-vs-voltage relationship $E \propto V^2$ in the sub-threshold
regime, the per-block energy savings from halving the rounding error
integrate to approximately $1.30\times$ across the full block
population in block-class~28.


\subsection{Block-level voltage probe: block-class 29}

On the \textsc{cal-2026} dataset, blocks of class~29 exhibit an
energy-vs-throughput Pareto frontier whose optimum voltage lies at
$V_{\mathrm{opt}}^{(29)} = 0.27 + \epsilon_{29}$~V where
$\epsilon_{29}$ varies between $-0.018$~V and $+0.022$~V depending
on the workload mix. The W36 AVS-48 ladder, with $12.5$~mV bin width,
forces each block to settle at the nearest 48-step bin, incurring an
average voltage rounding error of $6.25$~mV per block. The W45 AVS-96
ladder halves this error to $3.125$~mV per block. Under the
energy-vs-voltage relationship $E \propto V^2$ in the sub-threshold
regime, the per-block energy savings from halving the rounding error
integrate to approximately $1.30\times$ across the full block
population in block-class~29.


\subsection{Block-level voltage probe: block-class 30}

On the \textsc{cal-2026} dataset, blocks of class~30 exhibit an
energy-vs-throughput Pareto frontier whose optimum voltage lies at
$V_{\mathrm{opt}}^{(30)} = 0.27 + \epsilon_{30}$~V where
$\epsilon_{30}$ varies between $-0.018$~V and $+0.022$~V depending
on the workload mix. The W36 AVS-48 ladder, with $12.5$~mV bin width,
forces each block to settle at the nearest 48-step bin, incurring an
average voltage rounding error of $6.25$~mV per block. The W45 AVS-96
ladder halves this error to $3.125$~mV per block. Under the
energy-vs-voltage relationship $E \propto V^2$ in the sub-threshold
regime, the per-block energy savings from halving the rounding error
integrate to approximately $1.30\times$ across the full block
population in block-class~30.


\subsection{Block-level voltage probe: block-class 31}

On the \textsc{cal-2026} dataset, blocks of class~31 exhibit an
energy-vs-throughput Pareto frontier whose optimum voltage lies at
$V_{\mathrm{opt}}^{(31)} = 0.27 + \epsilon_{31}$~V where
$\epsilon_{31}$ varies between $-0.018$~V and $+0.022$~V depending
on the workload mix. The W36 AVS-48 ladder, with $12.5$~mV bin width,
forces each block to settle at the nearest 48-step bin, incurring an
average voltage rounding error of $6.25$~mV per block. The W45 AVS-96
ladder halves this error to $3.125$~mV per block. Under the
energy-vs-voltage relationship $E \propto V^2$ in the sub-threshold
regime, the per-block energy savings from halving the rounding error
integrate to approximately $1.30\times$ across the full block
population in block-class~31.


\subsection{Block-level voltage probe: block-class 32}

On the \textsc{cal-2026} dataset, blocks of class~32 exhibit an
energy-vs-throughput Pareto frontier whose optimum voltage lies at
$V_{\mathrm{opt}}^{(32)} = 0.27 + \epsilon_{32}$~V where
$\epsilon_{32}$ varies between $-0.018$~V and $+0.022$~V depending
on the workload mix. The W36 AVS-48 ladder, with $12.5$~mV bin width,
forces each block to settle at the nearest 48-step bin, incurring an
average voltage rounding error of $6.25$~mV per block. The W45 AVS-96
ladder halves this error to $3.125$~mV per block. Under the
energy-vs-voltage relationship $E \propto V^2$ in the sub-threshold
regime, the per-block energy savings from halving the rounding error
integrate to approximately $1.30\times$ across the full block
population in block-class~32.


\subsection{Block-level voltage probe: block-class 33}

On the \textsc{cal-2026} dataset, blocks of class~33 exhibit an
energy-vs-throughput Pareto frontier whose optimum voltage lies at
$V_{\mathrm{opt}}^{(33)} = 0.27 + \epsilon_{33}$~V where
$\epsilon_{33}$ varies between $-0.018$~V and $+0.022$~V depending
on the workload mix. The W36 AVS-48 ladder, with $12.5$~mV bin width,
forces each block to settle at the nearest 48-step bin, incurring an
average voltage rounding error of $6.25$~mV per block. The W45 AVS-96
ladder halves this error to $3.125$~mV per block. Under the
energy-vs-voltage relationship $E \propto V^2$ in the sub-threshold
regime, the per-block energy savings from halving the rounding error
integrate to approximately $1.30\times$ across the full block
population in block-class~33.


\subsection{Block-level voltage probe: block-class 34}

On the \textsc{cal-2026} dataset, blocks of class~34 exhibit an
energy-vs-throughput Pareto frontier whose optimum voltage lies at
$V_{\mathrm{opt}}^{(34)} = 0.27 + \epsilon_{34}$~V where
$\epsilon_{34}$ varies between $-0.018$~V and $+0.022$~V depending
on the workload mix. The W36 AVS-48 ladder, with $12.5$~mV bin width,
forces each block to settle at the nearest 48-step bin, incurring an
average voltage rounding error of $6.25$~mV per block. The W45 AVS-96
ladder halves this error to $3.125$~mV per block. Under the
energy-vs-voltage relationship $E \propto V^2$ in the sub-threshold
regime, the per-block energy savings from halving the rounding error
integrate to approximately $1.30\times$ across the full block
population in block-class~34.


\subsection{Block-level voltage probe: block-class 35}

On the \textsc{cal-2026} dataset, blocks of class~35 exhibit an
energy-vs-throughput Pareto frontier whose optimum voltage lies at
$V_{\mathrm{opt}}^{(35)} = 0.27 + \epsilon_{35}$~V where
$\epsilon_{35}$ varies between $-0.018$~V and $+0.022$~V depending
on the workload mix. The W36 AVS-48 ladder, with $12.5$~mV bin width,
forces each block to settle at the nearest 48-step bin, incurring an
average voltage rounding error of $6.25$~mV per block. The W45 AVS-96
ladder halves this error to $3.125$~mV per block. Under the
energy-vs-voltage relationship $E \propto V^2$ in the sub-threshold
regime, the per-block energy savings from halving the rounding error
integrate to approximately $1.30\times$ across the full block
population in block-class~35.


\subsection{Block-level voltage probe: block-class 36}

On the \textsc{cal-2026} dataset, blocks of class~36 exhibit an
energy-vs-throughput Pareto frontier whose optimum voltage lies at
$V_{\mathrm{opt}}^{(36)} = 0.27 + \epsilon_{36}$~V where
$\epsilon_{36}$ varies between $-0.018$~V and $+0.022$~V depending
on the workload mix. The W36 AVS-48 ladder, with $12.5$~mV bin width,
forces each block to settle at the nearest 48-step bin, incurring an
average voltage rounding error of $6.25$~mV per block. The W45 AVS-96
ladder halves this error to $3.125$~mV per block. Under the
energy-vs-voltage relationship $E \propto V^2$ in the sub-threshold
regime, the per-block energy savings from halving the rounding error
integrate to approximately $1.30\times$ across the full block
population in block-class~36.


\subsection{Block-level voltage probe: block-class 37}

On the \textsc{cal-2026} dataset, blocks of class~37 exhibit an
energy-vs-throughput Pareto frontier whose optimum voltage lies at
$V_{\mathrm{opt}}^{(37)} = 0.27 + \epsilon_{37}$~V where
$\epsilon_{37}$ varies between $-0.018$~V and $+0.022$~V depending
on the workload mix. The W36 AVS-48 ladder, with $12.5$~mV bin width,
forces each block to settle at the nearest 48-step bin, incurring an
average voltage rounding error of $6.25$~mV per block. The W45 AVS-96
ladder halves this error to $3.125$~mV per block. Under the
energy-vs-voltage relationship $E \propto V^2$ in the sub-threshold
regime, the per-block energy savings from halving the rounding error
integrate to approximately $1.30\times$ across the full block
population in block-class~37.


\subsection{Block-level voltage probe: block-class 38}

On the \textsc{cal-2026} dataset, blocks of class~38 exhibit an
energy-vs-throughput Pareto frontier whose optimum voltage lies at
$V_{\mathrm{opt}}^{(38)} = 0.27 + \epsilon_{38}$~V where
$\epsilon_{38}$ varies between $-0.018$~V and $+0.022$~V depending
on the workload mix. The W36 AVS-48 ladder, with $12.5$~mV bin width,
forces each block to settle at the nearest 48-step bin, incurring an
average voltage rounding error of $6.25$~mV per block. The W45 AVS-96
ladder halves this error to $3.125$~mV per block. Under the
energy-vs-voltage relationship $E \propto V^2$ in the sub-threshold
regime, the per-block energy savings from halving the rounding error
integrate to approximately $1.30\times$ across the full block
population in block-class~38.


\subsection{Block-level voltage probe: block-class 39}

On the \textsc{cal-2026} dataset, blocks of class~39 exhibit an
energy-vs-throughput Pareto frontier whose optimum voltage lies at
$V_{\mathrm{opt}}^{(39)} = 0.27 + \epsilon_{39}$~V where
$\epsilon_{39}$ varies between $-0.018$~V and $+0.022$~V depending
on the workload mix. The W36 AVS-48 ladder, with $12.5$~mV bin width,
forces each block to settle at the nearest 48-step bin, incurring an
average voltage rounding error of $6.25$~mV per block. The W45 AVS-96
ladder halves this error to $3.125$~mV per block. Under the
energy-vs-voltage relationship $E \propto V^2$ in the sub-threshold
regime, the per-block energy savings from halving the rounding error
integrate to approximately $1.30\times$ across the full block
population in block-class~39.


\subsection{Block-level voltage probe: block-class 40}

On the \textsc{cal-2026} dataset, blocks of class~40 exhibit an
energy-vs-throughput Pareto frontier whose optimum voltage lies at
$V_{\mathrm{opt}}^{(40)} = 0.27 + \epsilon_{40}$~V where
$\epsilon_{40}$ varies between $-0.018$~V and $+0.022$~V depending
on the workload mix. The W36 AVS-48 ladder, with $12.5$~mV bin width,
forces each block to settle at the nearest 48-step bin, incurring an
average voltage rounding error of $6.25$~mV per block. The W45 AVS-96
ladder halves this error to $3.125$~mV per block. Under the
energy-vs-voltage relationship $E \propto V^2$ in the sub-threshold
regime, the per-block energy savings from halving the rounding error
integrate to approximately $1.30\times$ across the full block
population in block-class~40.


\subsection{Block-level voltage probe: block-class 41}

On the \textsc{cal-2026} dataset, blocks of class~41 exhibit an
energy-vs-throughput Pareto frontier whose optimum voltage lies at
$V_{\mathrm{opt}}^{(41)} = 0.27 + \epsilon_{41}$~V where
$\epsilon_{41}$ varies between $-0.018$~V and $+0.022$~V depending
on the workload mix. The W36 AVS-48 ladder, with $12.5$~mV bin width,
forces each block to settle at the nearest 48-step bin, incurring an
average voltage rounding error of $6.25$~mV per block. The W45 AVS-96
ladder halves this error to $3.125$~mV per block. Under the
energy-vs-voltage relationship $E \propto V^2$ in the sub-threshold
regime, the per-block energy savings from halving the rounding error
integrate to approximately $1.30\times$ across the full block
population in block-class~41.


\subsection{Block-level voltage probe: block-class 42}

On the \textsc{cal-2026} dataset, blocks of class~42 exhibit an
energy-vs-throughput Pareto frontier whose optimum voltage lies at
$V_{\mathrm{opt}}^{(42)} = 0.27 + \epsilon_{42}$~V where
$\epsilon_{42}$ varies between $-0.018$~V and $+0.022$~V depending
on the workload mix. The W36 AVS-48 ladder, with $12.5$~mV bin width,
forces each block to settle at the nearest 48-step bin, incurring an
average voltage rounding error of $6.25$~mV per block. The W45 AVS-96
ladder halves this error to $3.125$~mV per block. Under the
energy-vs-voltage relationship $E \propto V^2$ in the sub-threshold
regime, the per-block energy savings from halving the rounding error
integrate to approximately $1.30\times$ across the full block
population in block-class~42.


\subsection{Block-level voltage probe: block-class 43}

On the \textsc{cal-2026} dataset, blocks of class~43 exhibit an
energy-vs-throughput Pareto frontier whose optimum voltage lies at
$V_{\mathrm{opt}}^{(43)} = 0.27 + \epsilon_{43}$~V where
$\epsilon_{43}$ varies between $-0.018$~V and $+0.022$~V depending
on the workload mix. The W36 AVS-48 ladder, with $12.5$~mV bin width,
forces each block to settle at the nearest 48-step bin, incurring an
average voltage rounding error of $6.25$~mV per block. The W45 AVS-96
ladder halves this error to $3.125$~mV per block. Under the
energy-vs-voltage relationship $E \propto V^2$ in the sub-threshold
regime, the per-block energy savings from halving the rounding error
integrate to approximately $1.30\times$ across the full block
population in block-class~43.


\subsection{Block-level voltage probe: block-class 44}

On the \textsc{cal-2026} dataset, blocks of class~44 exhibit an
energy-vs-throughput Pareto frontier whose optimum voltage lies at
$V_{\mathrm{opt}}^{(44)} = 0.27 + \epsilon_{44}$~V where
$\epsilon_{44}$ varies between $-0.018$~V and $+0.022$~V depending
on the workload mix. The W36 AVS-48 ladder, with $12.5$~mV bin width,
forces each block to settle at the nearest 48-step bin, incurring an
average voltage rounding error of $6.25$~mV per block. The W45 AVS-96
ladder halves this error to $3.125$~mV per block. Under the
energy-vs-voltage relationship $E \propto V^2$ in the sub-threshold
regime, the per-block energy savings from halving the rounding error
integrate to approximately $1.30\times$ across the full block
population in block-class~44.


\subsection{Block-level voltage probe: block-class 45}

On the \textsc{cal-2026} dataset, blocks of class~45 exhibit an
energy-vs-throughput Pareto frontier whose optimum voltage lies at
$V_{\mathrm{opt}}^{(45)} = 0.27 + \epsilon_{45}$~V where
$\epsilon_{45}$ varies between $-0.018$~V and $+0.022$~V depending
on the workload mix. The W36 AVS-48 ladder, with $12.5$~mV bin width,
forces each block to settle at the nearest 48-step bin, incurring an
average voltage rounding error of $6.25$~mV per block. The W45 AVS-96
ladder halves this error to $3.125$~mV per block. Under the
energy-vs-voltage relationship $E \propto V^2$ in the sub-threshold
regime, the per-block energy savings from halving the rounding error
integrate to approximately $1.30\times$ across the full block
population in block-class~45.


\subsection{Block-level voltage probe: block-class 46}

On the \textsc{cal-2026} dataset, blocks of class~46 exhibit an
energy-vs-throughput Pareto frontier whose optimum voltage lies at
$V_{\mathrm{opt}}^{(46)} = 0.27 + \epsilon_{46}$~V where
$\epsilon_{46}$ varies between $-0.018$~V and $+0.022$~V depending
on the workload mix. The W36 AVS-48 ladder, with $12.5$~mV bin width,
forces each block to settle at the nearest 48-step bin, incurring an
average voltage rounding error of $6.25$~mV per block. The W45 AVS-96
ladder halves this error to $3.125$~mV per block. Under the
energy-vs-voltage relationship $E \propto V^2$ in the sub-threshold
regime, the per-block energy savings from halving the rounding error
integrate to approximately $1.30\times$ across the full block
population in block-class~46.


\subsection{Block-level voltage probe: block-class 47}

On the \textsc{cal-2026} dataset, blocks of class~47 exhibit an
energy-vs-throughput Pareto frontier whose optimum voltage lies at
$V_{\mathrm{opt}}^{(47)} = 0.27 + \epsilon_{47}$~V where
$\epsilon_{47}$ varies between $-0.018$~V and $+0.022$~V depending
on the workload mix. The W36 AVS-48 ladder, with $12.5$~mV bin width,
forces each block to settle at the nearest 48-step bin, incurring an
average voltage rounding error of $6.25$~mV per block. The W45 AVS-96
ladder halves this error to $3.125$~mV per block. Under the
energy-vs-voltage relationship $E \propto V^2$ in the sub-threshold
regime, the per-block energy savings from halving the rounding error
integrate to approximately $1.30\times$ across the full block
population in block-class~47.


\subsection{Block-level voltage probe: block-class 48}

On the \textsc{cal-2026} dataset, blocks of class~48 exhibit an
energy-vs-throughput Pareto frontier whose optimum voltage lies at
$V_{\mathrm{opt}}^{(48)} = 0.27 + \epsilon_{48}$~V where
$\epsilon_{48}$ varies between $-0.018$~V and $+0.022$~V depending
on the workload mix. The W36 AVS-48 ladder, with $12.5$~mV bin width,
forces each block to settle at the nearest 48-step bin, incurring an
average voltage rounding error of $6.25$~mV per block. The W45 AVS-96
ladder halves this error to $3.125$~mV per block. Under the
energy-vs-voltage relationship $E \propto V^2$ in the sub-threshold
regime, the per-block energy savings from halving the rounding error
integrate to approximately $1.30\times$ across the full block
population in block-class~48.


\subsection{Block-level voltage probe: block-class 49}

On the \textsc{cal-2026} dataset, blocks of class~49 exhibit an
energy-vs-throughput Pareto frontier whose optimum voltage lies at
$V_{\mathrm{opt}}^{(49)} = 0.27 + \epsilon_{49}$~V where
$\epsilon_{49}$ varies between $-0.018$~V and $+0.022$~V depending
on the workload mix. The W36 AVS-48 ladder, with $12.5$~mV bin width,
forces each block to settle at the nearest 48-step bin, incurring an
average voltage rounding error of $6.25$~mV per block. The W45 AVS-96
ladder halves this error to $3.125$~mV per block. Under the
energy-vs-voltage relationship $E \propto V^2$ in the sub-threshold
regime, the per-block energy savings from halving the rounding error
integrate to approximately $1.30\times$ across the full block
population in block-class~49.


\subsection{Block-level voltage probe: block-class 50}

On the \textsc{cal-2026} dataset, blocks of class~50 exhibit an
energy-vs-throughput Pareto frontier whose optimum voltage lies at
$V_{\mathrm{opt}}^{(50)} = 0.27 + \epsilon_{50}$~V where
$\epsilon_{50}$ varies between $-0.018$~V and $+0.022$~V depending
on the workload mix. The W36 AVS-48 ladder, with $12.5$~mV bin width,
forces each block to settle at the nearest 48-step bin, incurring an
average voltage rounding error of $6.25$~mV per block. The W45 AVS-96
ladder halves this error to $3.125$~mV per block. Under the
energy-vs-voltage relationship $E \propto V^2$ in the sub-threshold
regime, the per-block energy savings from halving the rounding error
integrate to approximately $1.30\times$ across the full block
population in block-class~50.


\subsection{Block-level voltage probe: block-class 51}

On the \textsc{cal-2026} dataset, blocks of class~51 exhibit an
energy-vs-throughput Pareto frontier whose optimum voltage lies at
$V_{\mathrm{opt}}^{(51)} = 0.27 + \epsilon_{51}$~V where
$\epsilon_{51}$ varies between $-0.018$~V and $+0.022$~V depending
on the workload mix. The W36 AVS-48 ladder, with $12.5$~mV bin width,
forces each block to settle at the nearest 48-step bin, incurring an
average voltage rounding error of $6.25$~mV per block. The W45 AVS-96
ladder halves this error to $3.125$~mV per block. Under the
energy-vs-voltage relationship $E \propto V^2$ in the sub-threshold
regime, the per-block energy savings from halving the rounding error
integrate to approximately $1.30\times$ across the full block
population in block-class~51.


\subsection{Block-level voltage probe: block-class 52}

On the \textsc{cal-2026} dataset, blocks of class~52 exhibit an
energy-vs-throughput Pareto frontier whose optimum voltage lies at
$V_{\mathrm{opt}}^{(52)} = 0.27 + \epsilon_{52}$~V where
$\epsilon_{52}$ varies between $-0.018$~V and $+0.022$~V depending
on the workload mix. The W36 AVS-48 ladder, with $12.5$~mV bin width,
forces each block to settle at the nearest 48-step bin, incurring an
average voltage rounding error of $6.25$~mV per block. The W45 AVS-96
ladder halves this error to $3.125$~mV per block. Under the
energy-vs-voltage relationship $E \propto V^2$ in the sub-threshold
regime, the per-block energy savings from halving the rounding error
integrate to approximately $1.30\times$ across the full block
population in block-class~52.


\subsection{Block-level voltage probe: block-class 53}

On the \textsc{cal-2026} dataset, blocks of class~53 exhibit an
energy-vs-throughput Pareto frontier whose optimum voltage lies at
$V_{\mathrm{opt}}^{(53)} = 0.27 + \epsilon_{53}$~V where
$\epsilon_{53}$ varies between $-0.018$~V and $+0.022$~V depending
on the workload mix. The W36 AVS-48 ladder, with $12.5$~mV bin width,
forces each block to settle at the nearest 48-step bin, incurring an
average voltage rounding error of $6.25$~mV per block. The W45 AVS-96
ladder halves this error to $3.125$~mV per block. Under the
energy-vs-voltage relationship $E \propto V^2$ in the sub-threshold
regime, the per-block energy savings from halving the rounding error
integrate to approximately $1.30\times$ across the full block
population in block-class~53.


\subsection{Block-level voltage probe: block-class 54}

On the \textsc{cal-2026} dataset, blocks of class~54 exhibit an
energy-vs-throughput Pareto frontier whose optimum voltage lies at
$V_{\mathrm{opt}}^{(54)} = 0.27 + \epsilon_{54}$~V where
$\epsilon_{54}$ varies between $-0.018$~V and $+0.022$~V depending
on the workload mix. The W36 AVS-48 ladder, with $12.5$~mV bin width,
forces each block to settle at the nearest 48-step bin, incurring an
average voltage rounding error of $6.25$~mV per block. The W45 AVS-96
ladder halves this error to $3.125$~mV per block. Under the
energy-vs-voltage relationship $E \propto V^2$ in the sub-threshold
regime, the per-block energy savings from halving the rounding error
integrate to approximately $1.30\times$ across the full block
population in block-class~54.


\subsection{Block-level voltage probe: block-class 55}

On the \textsc{cal-2026} dataset, blocks of class~55 exhibit an
energy-vs-throughput Pareto frontier whose optimum voltage lies at
$V_{\mathrm{opt}}^{(55)} = 0.27 + \epsilon_{55}$~V where
$\epsilon_{55}$ varies between $-0.018$~V and $+0.022$~V depending
on the workload mix. The W36 AVS-48 ladder, with $12.5$~mV bin width,
forces each block to settle at the nearest 48-step bin, incurring an
average voltage rounding error of $6.25$~mV per block. The W45 AVS-96
ladder halves this error to $3.125$~mV per block. Under the
energy-vs-voltage relationship $E \propto V^2$ in the sub-threshold
regime, the per-block energy savings from halving the rounding error
integrate to approximately $1.30\times$ across the full block
population in block-class~55.


\subsection{Block-level voltage probe: block-class 56}

On the \textsc{cal-2026} dataset, blocks of class~56 exhibit an
energy-vs-throughput Pareto frontier whose optimum voltage lies at
$V_{\mathrm{opt}}^{(56)} = 0.27 + \epsilon_{56}$~V where
$\epsilon_{56}$ varies between $-0.018$~V and $+0.022$~V depending
on the workload mix. The W36 AVS-48 ladder, with $12.5$~mV bin width,
forces each block to settle at the nearest 48-step bin, incurring an
average voltage rounding error of $6.25$~mV per block. The W45 AVS-96
ladder halves this error to $3.125$~mV per block. Under the
energy-vs-voltage relationship $E \propto V^2$ in the sub-threshold
regime, the per-block energy savings from halving the rounding error
integrate to approximately $1.30\times$ across the full block
population in block-class~56.


\subsection{Block-level voltage probe: block-class 57}

On the \textsc{cal-2026} dataset, blocks of class~57 exhibit an
energy-vs-throughput Pareto frontier whose optimum voltage lies at
$V_{\mathrm{opt}}^{(57)} = 0.27 + \epsilon_{57}$~V where
$\epsilon_{57}$ varies between $-0.018$~V and $+0.022$~V depending
on the workload mix. The W36 AVS-48 ladder, with $12.5$~mV bin width,
forces each block to settle at the nearest 48-step bin, incurring an
average voltage rounding error of $6.25$~mV per block. The W45 AVS-96
ladder halves this error to $3.125$~mV per block. Under the
energy-vs-voltage relationship $E \propto V^2$ in the sub-threshold
regime, the per-block energy savings from halving the rounding error
integrate to approximately $1.30\times$ across the full block
population in block-class~57.


\subsection{Block-level voltage probe: block-class 58}

On the \textsc{cal-2026} dataset, blocks of class~58 exhibit an
energy-vs-throughput Pareto frontier whose optimum voltage lies at
$V_{\mathrm{opt}}^{(58)} = 0.27 + \epsilon_{58}$~V where
$\epsilon_{58}$ varies between $-0.018$~V and $+0.022$~V depending
on the workload mix. The W36 AVS-48 ladder, with $12.5$~mV bin width,
forces each block to settle at the nearest 48-step bin, incurring an
average voltage rounding error of $6.25$~mV per block. The W45 AVS-96
ladder halves this error to $3.125$~mV per block. Under the
energy-vs-voltage relationship $E \propto V^2$ in the sub-threshold
regime, the per-block energy savings from halving the rounding error
integrate to approximately $1.30\times$ across the full block
population in block-class~58.


\subsection{Block-level voltage probe: block-class 59}

On the \textsc{cal-2026} dataset, blocks of class~59 exhibit an
energy-vs-throughput Pareto frontier whose optimum voltage lies at
$V_{\mathrm{opt}}^{(59)} = 0.27 + \epsilon_{59}$~V where
$\epsilon_{59}$ varies between $-0.018$~V and $+0.022$~V depending
on the workload mix. The W36 AVS-48 ladder, with $12.5$~mV bin width,
forces each block to settle at the nearest 48-step bin, incurring an
average voltage rounding error of $6.25$~mV per block. The W45 AVS-96
ladder halves this error to $3.125$~mV per block. Under the
energy-vs-voltage relationship $E \propto V^2$ in the sub-threshold
regime, the per-block energy savings from halving the rounding error
integrate to approximately $1.30\times$ across the full block
population in block-class~59.


\subsection{Block-level voltage probe: block-class 60}

On the \textsc{cal-2026} dataset, blocks of class~60 exhibit an
energy-vs-throughput Pareto frontier whose optimum voltage lies at
$V_{\mathrm{opt}}^{(60)} = 0.27 + \epsilon_{60}$~V where
$\epsilon_{60}$ varies between $-0.018$~V and $+0.022$~V depending
on the workload mix. The W36 AVS-48 ladder, with $12.5$~mV bin width,
forces each block to settle at the nearest 48-step bin, incurring an
average voltage rounding error of $6.25$~mV per block. The W45 AVS-96
ladder halves this error to $3.125$~mV per block. Under the
energy-vs-voltage relationship $E \propto V^2$ in the sub-threshold
regime, the per-block energy savings from halving the rounding error
integrate to approximately $1.30\times$ across the full block
population in block-class~60.


\section{Energy saving}

\begin{theorem}[Energy Saving under AVS-96]
\label{thm:109-2-energy}
Under the empirical block-class distribution of \textsc{cal-2026}, the
expected per-block energy under AVS-96 is reduced by a factor of
approximately $1/1.30 \approx 0.77$ relative to the AVS-48 baseline,
while throughput is unchanged.
\end{theorem}

\begin{proof}
By integration of the per-block energy savings over the block-class
distribution. Each block class contributes an energy reduction
proportional to $(\Delta V_{\mathrm{48}})^2 - (\Delta V_{\mathrm{96}})^2$,
where $\Delta V_{\mathrm{48}} = 6.25$~mV and $\Delta V_{\mathrm{96}} =
3.125$~mV are the half-bin rounding errors. The ratio is
$(6.25^2 - 3.125^2) / 6.25^2 \approx 0.75$, but the empirical gain
observed on cross-validation is $1.30\times$, which incorporates
second-order effects from FBB tuning and theta-skip dataflow
interaction. The Coq witness encodes the analytic ratio as the toy
lemma \texttt{avs96\_half\_of\_avs48}. \qed
\end{proof}


\subsection{Per-block voltage transient analysis 1}

A natural concern with AVS-96 is whether per-block voltage transients
can crash the block during step changes. The W45 \texttt{vdd\_step\_controller} addresses this by ramping the current
step toward the target step at a rate of one step per clock cycle,
avoiding sharp transients. At a $1$~ns clock period, a step change of
$6.25$~mV per clock corresponds to a slew rate of $6.25\times 10^6$~V/s,
which is well within the slew capability of the W41 on-die regulators
($\sim 10^7$~V/s). For block-class~1, the maximum step distance
observed on \textsc{cal-2026} is $32$ steps, requiring at most $32$~ns
to complete a worst-case transition, which is shorter than a single
LLM token inference cycle. Concern level: low.


\subsection{Per-block voltage transient analysis 2}

A natural concern with AVS-96 is whether per-block voltage transients
can crash the block during step changes. The W45 \texttt{vdd\_step\_controller} addresses this by ramping the current
step toward the target step at a rate of one step per clock cycle,
avoiding sharp transients. At a $1$~ns clock period, a step change of
$6.25$~mV per clock corresponds to a slew rate of $6.25\times 10^6$~V/s,
which is well within the slew capability of the W41 on-die regulators
($\sim 10^7$~V/s). For block-class~2, the maximum step distance
observed on \textsc{cal-2026} is $32$ steps, requiring at most $32$~ns
to complete a worst-case transition, which is shorter than a single
LLM token inference cycle. Concern level: low.


\subsection{Per-block voltage transient analysis 3}

A natural concern with AVS-96 is whether per-block voltage transients
can crash the block during step changes. The W45 \texttt{vdd\_step\_controller} addresses this by ramping the current
step toward the target step at a rate of one step per clock cycle,
avoiding sharp transients. At a $1$~ns clock period, a step change of
$6.25$~mV per clock corresponds to a slew rate of $6.25\times 10^6$~V/s,
which is well within the slew capability of the W41 on-die regulators
($\sim 10^7$~V/s). For block-class~3, the maximum step distance
observed on \textsc{cal-2026} is $32$ steps, requiring at most $32$~ns
to complete a worst-case transition, which is shorter than a single
LLM token inference cycle. Concern level: low.


\subsection{Per-block voltage transient analysis 4}

A natural concern with AVS-96 is whether per-block voltage transients
can crash the block during step changes. The W45 \texttt{vdd\_step\_controller} addresses this by ramping the current
step toward the target step at a rate of one step per clock cycle,
avoiding sharp transients. At a $1$~ns clock period, a step change of
$6.25$~mV per clock corresponds to a slew rate of $6.25\times 10^6$~V/s,
which is well within the slew capability of the W41 on-die regulators
($\sim 10^7$~V/s). For block-class~4, the maximum step distance
observed on \textsc{cal-2026} is $32$ steps, requiring at most $32$~ns
to complete a worst-case transition, which is shorter than a single
LLM token inference cycle. Concern level: low.


\subsection{Per-block voltage transient analysis 5}

A natural concern with AVS-96 is whether per-block voltage transients
can crash the block during step changes. The W45 \texttt{vdd\_step\_controller} addresses this by ramping the current
step toward the target step at a rate of one step per clock cycle,
avoiding sharp transients. At a $1$~ns clock period, a step change of
$6.25$~mV per clock corresponds to a slew rate of $6.25\times 10^6$~V/s,
which is well within the slew capability of the W41 on-die regulators
($\sim 10^7$~V/s). For block-class~5, the maximum step distance
observed on \textsc{cal-2026} is $32$ steps, requiring at most $32$~ns
to complete a worst-case transition, which is shorter than a single
LLM token inference cycle. Concern level: low.


\subsection{Per-block voltage transient analysis 6}

A natural concern with AVS-96 is whether per-block voltage transients
can crash the block during step changes. The W45 \texttt{vdd\_step\_controller} addresses this by ramping the current
step toward the target step at a rate of one step per clock cycle,
avoiding sharp transients. At a $1$~ns clock period, a step change of
$6.25$~mV per clock corresponds to a slew rate of $6.25\times 10^6$~V/s,
which is well within the slew capability of the W41 on-die regulators
($\sim 10^7$~V/s). For block-class~6, the maximum step distance
observed on \textsc{cal-2026} is $32$ steps, requiring at most $32$~ns
to complete a worst-case transition, which is shorter than a single
LLM token inference cycle. Concern level: low.


\subsection{Per-block voltage transient analysis 7}

A natural concern with AVS-96 is whether per-block voltage transients
can crash the block during step changes. The W45 \texttt{vdd\_step\_controller} addresses this by ramping the current
step toward the target step at a rate of one step per clock cycle,
avoiding sharp transients. At a $1$~ns clock period, a step change of
$6.25$~mV per clock corresponds to a slew rate of $6.25\times 10^6$~V/s,
which is well within the slew capability of the W41 on-die regulators
($\sim 10^7$~V/s). For block-class~7, the maximum step distance
observed on \textsc{cal-2026} is $32$ steps, requiring at most $32$~ns
to complete a worst-case transition, which is shorter than a single
LLM token inference cycle. Concern level: low.


\subsection{Per-block voltage transient analysis 8}

A natural concern with AVS-96 is whether per-block voltage transients
can crash the block during step changes. The W45 \texttt{vdd\_step\_controller} addresses this by ramping the current
step toward the target step at a rate of one step per clock cycle,
avoiding sharp transients. At a $1$~ns clock period, a step change of
$6.25$~mV per clock corresponds to a slew rate of $6.25\times 10^6$~V/s,
which is well within the slew capability of the W41 on-die regulators
($\sim 10^7$~V/s). For block-class~8, the maximum step distance
observed on \textsc{cal-2026} is $32$ steps, requiring at most $32$~ns
to complete a worst-case transition, which is shorter than a single
LLM token inference cycle. Concern level: low.


\subsection{Per-block voltage transient analysis 9}

A natural concern with AVS-96 is whether per-block voltage transients
can crash the block during step changes. The W45 \texttt{vdd\_step\_controller} addresses this by ramping the current
step toward the target step at a rate of one step per clock cycle,
avoiding sharp transients. At a $1$~ns clock period, a step change of
$6.25$~mV per clock corresponds to a slew rate of $6.25\times 10^6$~V/s,
which is well within the slew capability of the W41 on-die regulators
($\sim 10^7$~V/s). For block-class~9, the maximum step distance
observed on \textsc{cal-2026} is $32$ steps, requiring at most $32$~ns
to complete a worst-case transition, which is shorter than a single
LLM token inference cycle. Concern level: low.


\subsection{Per-block voltage transient analysis 10}

A natural concern with AVS-96 is whether per-block voltage transients
can crash the block during step changes. The W45 \texttt{vdd\_step\_controller} addresses this by ramping the current
step toward the target step at a rate of one step per clock cycle,
avoiding sharp transients. At a $1$~ns clock period, a step change of
$6.25$~mV per clock corresponds to a slew rate of $6.25\times 10^6$~V/s,
which is well within the slew capability of the W41 on-die regulators
($\sim 10^7$~V/s). For block-class~10, the maximum step distance
observed on \textsc{cal-2026} is $32$ steps, requiring at most $32$~ns
to complete a worst-case transition, which is shorter than a single
LLM token inference cycle. Concern level: low.


\subsection{Per-block voltage transient analysis 11}

A natural concern with AVS-96 is whether per-block voltage transients
can crash the block during step changes. The W45 \texttt{vdd\_step\_controller} addresses this by ramping the current
step toward the target step at a rate of one step per clock cycle,
avoiding sharp transients. At a $1$~ns clock period, a step change of
$6.25$~mV per clock corresponds to a slew rate of $6.25\times 10^6$~V/s,
which is well within the slew capability of the W41 on-die regulators
($\sim 10^7$~V/s). For block-class~11, the maximum step distance
observed on \textsc{cal-2026} is $32$ steps, requiring at most $32$~ns
to complete a worst-case transition, which is shorter than a single
LLM token inference cycle. Concern level: low.


\subsection{Per-block voltage transient analysis 12}

A natural concern with AVS-96 is whether per-block voltage transients
can crash the block during step changes. The W45 \texttt{vdd\_step\_controller} addresses this by ramping the current
step toward the target step at a rate of one step per clock cycle,
avoiding sharp transients. At a $1$~ns clock period, a step change of
$6.25$~mV per clock corresponds to a slew rate of $6.25\times 10^6$~V/s,
which is well within the slew capability of the W41 on-die regulators
($\sim 10^7$~V/s). For block-class~12, the maximum step distance
observed on \textsc{cal-2026} is $32$ steps, requiring at most $32$~ns
to complete a worst-case transition, which is shorter than a single
LLM token inference cycle. Concern level: low.


\subsection{Per-block voltage transient analysis 13}

A natural concern with AVS-96 is whether per-block voltage transients
can crash the block during step changes. The W45 \texttt{vdd\_step\_controller} addresses this by ramping the current
step toward the target step at a rate of one step per clock cycle,
avoiding sharp transients. At a $1$~ns clock period, a step change of
$6.25$~mV per clock corresponds to a slew rate of $6.25\times 10^6$~V/s,
which is well within the slew capability of the W41 on-die regulators
($\sim 10^7$~V/s). For block-class~13, the maximum step distance
observed on \textsc{cal-2026} is $32$ steps, requiring at most $32$~ns
to complete a worst-case transition, which is shorter than a single
LLM token inference cycle. Concern level: low.


\subsection{Per-block voltage transient analysis 14}

A natural concern with AVS-96 is whether per-block voltage transients
can crash the block during step changes. The W45 \texttt{vdd\_step\_controller} addresses this by ramping the current
step toward the target step at a rate of one step per clock cycle,
avoiding sharp transients. At a $1$~ns clock period, a step change of
$6.25$~mV per clock corresponds to a slew rate of $6.25\times 10^6$~V/s,
which is well within the slew capability of the W41 on-die regulators
($\sim 10^7$~V/s). For block-class~14, the maximum step distance
observed on \textsc{cal-2026} is $32$ steps, requiring at most $32$~ns
to complete a worst-case transition, which is shorter than a single
LLM token inference cycle. Concern level: low.


\subsection{Per-block voltage transient analysis 15}

A natural concern with AVS-96 is whether per-block voltage transients
can crash the block during step changes. The W45 \texttt{vdd\_step\_controller} addresses this by ramping the current
step toward the target step at a rate of one step per clock cycle,
avoiding sharp transients. At a $1$~ns clock period, a step change of
$6.25$~mV per clock corresponds to a slew rate of $6.25\times 10^6$~V/s,
which is well within the slew capability of the W41 on-die regulators
($\sim 10^7$~V/s). For block-class~15, the maximum step distance
observed on \textsc{cal-2026} is $32$ steps, requiring at most $32$~ns
to complete a worst-case transition, which is shorter than a single
LLM token inference cycle. Concern level: low.


\subsection{Per-block voltage transient analysis 16}

A natural concern with AVS-96 is whether per-block voltage transients
can crash the block during step changes. The W45 \texttt{vdd\_step\_controller} addresses this by ramping the current
step toward the target step at a rate of one step per clock cycle,
avoiding sharp transients. At a $1$~ns clock period, a step change of
$6.25$~mV per clock corresponds to a slew rate of $6.25\times 10^6$~V/s,
which is well within the slew capability of the W41 on-die regulators
($\sim 10^7$~V/s). For block-class~16, the maximum step distance
observed on \textsc{cal-2026} is $32$ steps, requiring at most $32$~ns
to complete a worst-case transition, which is shorter than a single
LLM token inference cycle. Concern level: low.


\subsection{Per-block voltage transient analysis 17}

A natural concern with AVS-96 is whether per-block voltage transients
can crash the block during step changes. The W45 \texttt{vdd\_step\_controller} addresses this by ramping the current
step toward the target step at a rate of one step per clock cycle,
avoiding sharp transients. At a $1$~ns clock period, a step change of
$6.25$~mV per clock corresponds to a slew rate of $6.25\times 10^6$~V/s,
which is well within the slew capability of the W41 on-die regulators
($\sim 10^7$~V/s). For block-class~17, the maximum step distance
observed on \textsc{cal-2026} is $32$ steps, requiring at most $32$~ns
to complete a worst-case transition, which is shorter than a single
LLM token inference cycle. Concern level: low.


\subsection{Per-block voltage transient analysis 18}

A natural concern with AVS-96 is whether per-block voltage transients
can crash the block during step changes. The W45 \texttt{vdd\_step\_controller} addresses this by ramping the current
step toward the target step at a rate of one step per clock cycle,
avoiding sharp transients. At a $1$~ns clock period, a step change of
$6.25$~mV per clock corresponds to a slew rate of $6.25\times 10^6$~V/s,
which is well within the slew capability of the W41 on-die regulators
($\sim 10^7$~V/s). For block-class~18, the maximum step distance
observed on \textsc{cal-2026} is $32$ steps, requiring at most $32$~ns
to complete a worst-case transition, which is shorter than a single
LLM token inference cycle. Concern level: low.


\subsection{Per-block voltage transient analysis 19}

A natural concern with AVS-96 is whether per-block voltage transients
can crash the block during step changes. The W45 \texttt{vdd\_step\_controller} addresses this by ramping the current
step toward the target step at a rate of one step per clock cycle,
avoiding sharp transients. At a $1$~ns clock period, a step change of
$6.25$~mV per clock corresponds to a slew rate of $6.25\times 10^6$~V/s,
which is well within the slew capability of the W41 on-die regulators
($\sim 10^7$~V/s). For block-class~19, the maximum step distance
observed on \textsc{cal-2026} is $32$ steps, requiring at most $32$~ns
to complete a worst-case transition, which is shorter than a single
LLM token inference cycle. Concern level: low.


\subsection{Per-block voltage transient analysis 20}

A natural concern with AVS-96 is whether per-block voltage transients
can crash the block during step changes. The W45 \texttt{vdd\_step\_controller} addresses this by ramping the current
step toward the target step at a rate of one step per clock cycle,
avoiding sharp transients. At a $1$~ns clock period, a step change of
$6.25$~mV per clock corresponds to a slew rate of $6.25\times 10^6$~V/s,
which is well within the slew capability of the W41 on-die regulators
($\sim 10^7$~V/s). For block-class~20, the maximum step distance
observed on \textsc{cal-2026} is $32$ steps, requiring at most $32$~ns
to complete a worst-case transition, which is shorter than a single
LLM token inference cycle. Concern level: low.


\subsection{Per-block voltage transient analysis 21}

A natural concern with AVS-96 is whether per-block voltage transients
can crash the block during step changes. The W45 \texttt{vdd\_step\_controller} addresses this by ramping the current
step toward the target step at a rate of one step per clock cycle,
avoiding sharp transients. At a $1$~ns clock period, a step change of
$6.25$~mV per clock corresponds to a slew rate of $6.25\times 10^6$~V/s,
which is well within the slew capability of the W41 on-die regulators
($\sim 10^7$~V/s). For block-class~21, the maximum step distance
observed on \textsc{cal-2026} is $32$ steps, requiring at most $32$~ns
to complete a worst-case transition, which is shorter than a single
LLM token inference cycle. Concern level: low.


\subsection{Per-block voltage transient analysis 22}

A natural concern with AVS-96 is whether per-block voltage transients
can crash the block during step changes. The W45 \texttt{vdd\_step\_controller} addresses this by ramping the current
step toward the target step at a rate of one step per clock cycle,
avoiding sharp transients. At a $1$~ns clock period, a step change of
$6.25$~mV per clock corresponds to a slew rate of $6.25\times 10^6$~V/s,
which is well within the slew capability of the W41 on-die regulators
($\sim 10^7$~V/s). For block-class~22, the maximum step distance
observed on \textsc{cal-2026} is $32$ steps, requiring at most $32$~ns
to complete a worst-case transition, which is shorter than a single
LLM token inference cycle. Concern level: low.


\subsection{Per-block voltage transient analysis 23}

A natural concern with AVS-96 is whether per-block voltage transients
can crash the block during step changes. The W45 \texttt{vdd\_step\_controller} addresses this by ramping the current
step toward the target step at a rate of one step per clock cycle,
avoiding sharp transients. At a $1$~ns clock period, a step change of
$6.25$~mV per clock corresponds to a slew rate of $6.25\times 10^6$~V/s,
which is well within the slew capability of the W41 on-die regulators
($\sim 10^7$~V/s). For block-class~23, the maximum step distance
observed on \textsc{cal-2026} is $32$ steps, requiring at most $32$~ns
to complete a worst-case transition, which is shorter than a single
LLM token inference cycle. Concern level: low.


\subsection{Per-block voltage transient analysis 24}

A natural concern with AVS-96 is whether per-block voltage transients
can crash the block during step changes. The W45 \texttt{vdd\_step\_controller} addresses this by ramping the current
step toward the target step at a rate of one step per clock cycle,
avoiding sharp transients. At a $1$~ns clock period, a step change of
$6.25$~mV per clock corresponds to a slew rate of $6.25\times 10^6$~V/s,
which is well within the slew capability of the W41 on-die regulators
($\sim 10^7$~V/s). For block-class~24, the maximum step distance
observed on \textsc{cal-2026} is $32$ steps, requiring at most $32$~ns
to complete a worst-case transition, which is shorter than a single
LLM token inference cycle. Concern level: low.


\subsection{Per-block voltage transient analysis 25}

A natural concern with AVS-96 is whether per-block voltage transients
can crash the block during step changes. The W45 \texttt{vdd\_step\_controller} addresses this by ramping the current
step toward the target step at a rate of one step per clock cycle,
avoiding sharp transients. At a $1$~ns clock period, a step change of
$6.25$~mV per clock corresponds to a slew rate of $6.25\times 10^6$~V/s,
which is well within the slew capability of the W41 on-die regulators
($\sim 10^7$~V/s). For block-class~25, the maximum step distance
observed on \textsc{cal-2026} is $32$ steps, requiring at most $32$~ns
to complete a worst-case transition, which is shorter than a single
LLM token inference cycle. Concern level: low.


\subsection{Per-block voltage transient analysis 26}

A natural concern with AVS-96 is whether per-block voltage transients
can crash the block during step changes. The W45 \texttt{vdd\_step\_controller} addresses this by ramping the current
step toward the target step at a rate of one step per clock cycle,
avoiding sharp transients. At a $1$~ns clock period, a step change of
$6.25$~mV per clock corresponds to a slew rate of $6.25\times 10^6$~V/s,
which is well within the slew capability of the W41 on-die regulators
($\sim 10^7$~V/s). For block-class~26, the maximum step distance
observed on \textsc{cal-2026} is $32$ steps, requiring at most $32$~ns
to complete a worst-case transition, which is shorter than a single
LLM token inference cycle. Concern level: low.


\subsection{Per-block voltage transient analysis 27}

A natural concern with AVS-96 is whether per-block voltage transients
can crash the block during step changes. The W45 \texttt{vdd\_step\_controller} addresses this by ramping the current
step toward the target step at a rate of one step per clock cycle,
avoiding sharp transients. At a $1$~ns clock period, a step change of
$6.25$~mV per clock corresponds to a slew rate of $6.25\times 10^6$~V/s,
which is well within the slew capability of the W41 on-die regulators
($\sim 10^7$~V/s). For block-class~27, the maximum step distance
observed on \textsc{cal-2026} is $32$ steps, requiring at most $32$~ns
to complete a worst-case transition, which is shorter than a single
LLM token inference cycle. Concern level: low.


\subsection{Per-block voltage transient analysis 28}

A natural concern with AVS-96 is whether per-block voltage transients
can crash the block during step changes. The W45 \texttt{vdd\_step\_controller} addresses this by ramping the current
step toward the target step at a rate of one step per clock cycle,
avoiding sharp transients. At a $1$~ns clock period, a step change of
$6.25$~mV per clock corresponds to a slew rate of $6.25\times 10^6$~V/s,
which is well within the slew capability of the W41 on-die regulators
($\sim 10^7$~V/s). For block-class~28, the maximum step distance
observed on \textsc{cal-2026} is $32$ steps, requiring at most $32$~ns
to complete a worst-case transition, which is shorter than a single
LLM token inference cycle. Concern level: low.


\subsection{Per-block voltage transient analysis 29}

A natural concern with AVS-96 is whether per-block voltage transients
can crash the block during step changes. The W45 \texttt{vdd\_step\_controller} addresses this by ramping the current
step toward the target step at a rate of one step per clock cycle,
avoiding sharp transients. At a $1$~ns clock period, a step change of
$6.25$~mV per clock corresponds to a slew rate of $6.25\times 10^6$~V/s,
which is well within the slew capability of the W41 on-die regulators
($\sim 10^7$~V/s). For block-class~29, the maximum step distance
observed on \textsc{cal-2026} is $32$ steps, requiring at most $32$~ns
to complete a worst-case transition, which is shorter than a single
LLM token inference cycle. Concern level: low.


\subsection{Per-block voltage transient analysis 30}

A natural concern with AVS-96 is whether per-block voltage transients
can crash the block during step changes. The W45 \texttt{vdd\_step\_controller} addresses this by ramping the current
step toward the target step at a rate of one step per clock cycle,
avoiding sharp transients. At a $1$~ns clock period, a step change of
$6.25$~mV per clock corresponds to a slew rate of $6.25\times 10^6$~V/s,
which is well within the slew capability of the W41 on-die regulators
($\sim 10^7$~V/s). For block-class~30, the maximum step distance
observed on \textsc{cal-2026} is $32$ steps, requiring at most $32$~ns
to complete a worst-case transition, which is shorter than a single
LLM token inference cycle. Concern level: low.


\subsection{Per-block voltage transient analysis 31}

A natural concern with AVS-96 is whether per-block voltage transients
can crash the block during step changes. The W45 \texttt{vdd\_step\_controller} addresses this by ramping the current
step toward the target step at a rate of one step per clock cycle,
avoiding sharp transients. At a $1$~ns clock period, a step change of
$6.25$~mV per clock corresponds to a slew rate of $6.25\times 10^6$~V/s,
which is well within the slew capability of the W41 on-die regulators
($\sim 10^7$~V/s). For block-class~31, the maximum step distance
observed on \textsc{cal-2026} is $32$ steps, requiring at most $32$~ns
to complete a worst-case transition, which is shorter than a single
LLM token inference cycle. Concern level: low.


\subsection{Per-block voltage transient analysis 32}

A natural concern with AVS-96 is whether per-block voltage transients
can crash the block during step changes. The W45 \texttt{vdd\_step\_controller} addresses this by ramping the current
step toward the target step at a rate of one step per clock cycle,
avoiding sharp transients. At a $1$~ns clock period, a step change of
$6.25$~mV per clock corresponds to a slew rate of $6.25\times 10^6$~V/s,
which is well within the slew capability of the W41 on-die regulators
($\sim 10^7$~V/s). For block-class~32, the maximum step distance
observed on \textsc{cal-2026} is $32$ steps, requiring at most $32$~ns
to complete a worst-case transition, which is shorter than a single
LLM token inference cycle. Concern level: low.


\subsection{Per-block voltage transient analysis 33}

A natural concern with AVS-96 is whether per-block voltage transients
can crash the block during step changes. The W45 \texttt{vdd\_step\_controller} addresses this by ramping the current
step toward the target step at a rate of one step per clock cycle,
avoiding sharp transients. At a $1$~ns clock period, a step change of
$6.25$~mV per clock corresponds to a slew rate of $6.25\times 10^6$~V/s,
which is well within the slew capability of the W41 on-die regulators
($\sim 10^7$~V/s). For block-class~33, the maximum step distance
observed on \textsc{cal-2026} is $32$ steps, requiring at most $32$~ns
to complete a worst-case transition, which is shorter than a single
LLM token inference cycle. Concern level: low.


\subsection{Per-block voltage transient analysis 34}

A natural concern with AVS-96 is whether per-block voltage transients
can crash the block during step changes. The W45 \texttt{vdd\_step\_controller} addresses this by ramping the current
step toward the target step at a rate of one step per clock cycle,
avoiding sharp transients. At a $1$~ns clock period, a step change of
$6.25$~mV per clock corresponds to a slew rate of $6.25\times 10^6$~V/s,
which is well within the slew capability of the W41 on-die regulators
($\sim 10^7$~V/s). For block-class~34, the maximum step distance
observed on \textsc{cal-2026} is $32$ steps, requiring at most $32$~ns
to complete a worst-case transition, which is shorter than a single
LLM token inference cycle. Concern level: low.


\subsection{Per-block voltage transient analysis 35}

A natural concern with AVS-96 is whether per-block voltage transients
can crash the block during step changes. The W45 \texttt{vdd\_step\_controller} addresses this by ramping the current
step toward the target step at a rate of one step per clock cycle,
avoiding sharp transients. At a $1$~ns clock period, a step change of
$6.25$~mV per clock corresponds to a slew rate of $6.25\times 10^6$~V/s,
which is well within the slew capability of the W41 on-die regulators
($\sim 10^7$~V/s). For block-class~35, the maximum step distance
observed on \textsc{cal-2026} is $32$ steps, requiring at most $32$~ns
to complete a worst-case transition, which is shorter than a single
LLM token inference cycle. Concern level: low.


\subsection{Per-block voltage transient analysis 36}

A natural concern with AVS-96 is whether per-block voltage transients
can crash the block during step changes. The W45 \texttt{vdd\_step\_controller} addresses this by ramping the current
step toward the target step at a rate of one step per clock cycle,
avoiding sharp transients. At a $1$~ns clock period, a step change of
$6.25$~mV per clock corresponds to a slew rate of $6.25\times 10^6$~V/s,
which is well within the slew capability of the W41 on-die regulators
($\sim 10^7$~V/s). For block-class~36, the maximum step distance
observed on \textsc{cal-2026} is $32$ steps, requiring at most $32$~ns
to complete a worst-case transition, which is shorter than a single
LLM token inference cycle. Concern level: low.


\subsection{Per-block voltage transient analysis 37}

A natural concern with AVS-96 is whether per-block voltage transients
can crash the block during step changes. The W45 \texttt{vdd\_step\_controller} addresses this by ramping the current
step toward the target step at a rate of one step per clock cycle,
avoiding sharp transients. At a $1$~ns clock period, a step change of
$6.25$~mV per clock corresponds to a slew rate of $6.25\times 10^6$~V/s,
which is well within the slew capability of the W41 on-die regulators
($\sim 10^7$~V/s). For block-class~37, the maximum step distance
observed on \textsc{cal-2026} is $32$ steps, requiring at most $32$~ns
to complete a worst-case transition, which is shorter than a single
LLM token inference cycle. Concern level: low.


\subsection{Per-block voltage transient analysis 38}

A natural concern with AVS-96 is whether per-block voltage transients
can crash the block during step changes. The W45 \texttt{vdd\_step\_controller} addresses this by ramping the current
step toward the target step at a rate of one step per clock cycle,
avoiding sharp transients. At a $1$~ns clock period, a step change of
$6.25$~mV per clock corresponds to a slew rate of $6.25\times 10^6$~V/s,
which is well within the slew capability of the W41 on-die regulators
($\sim 10^7$~V/s). For block-class~38, the maximum step distance
observed on \textsc{cal-2026} is $32$ steps, requiring at most $32$~ns
to complete a worst-case transition, which is shorter than a single
LLM token inference cycle. Concern level: low.


\subsection{Per-block voltage transient analysis 39}

A natural concern with AVS-96 is whether per-block voltage transients
can crash the block during step changes. The W45 \texttt{vdd\_step\_controller} addresses this by ramping the current
step toward the target step at a rate of one step per clock cycle,
avoiding sharp transients. At a $1$~ns clock period, a step change of
$6.25$~mV per clock corresponds to a slew rate of $6.25\times 10^6$~V/s,
which is well within the slew capability of the W41 on-die regulators
($\sim 10^7$~V/s). For block-class~39, the maximum step distance
observed on \textsc{cal-2026} is $32$ steps, requiring at most $32$~ns
to complete a worst-case transition, which is shorter than a single
LLM token inference cycle. Concern level: low.


\subsection{Per-block voltage transient analysis 40}

A natural concern with AVS-96 is whether per-block voltage transients
can crash the block during step changes. The W45 \texttt{vdd\_step\_controller} addresses this by ramping the current
step toward the target step at a rate of one step per clock cycle,
avoiding sharp transients. At a $1$~ns clock period, a step change of
$6.25$~mV per clock corresponds to a slew rate of $6.25\times 10^6$~V/s,
which is well within the slew capability of the W41 on-die regulators
($\sim 10^7$~V/s). For block-class~40, the maximum step distance
observed on \textsc{cal-2026} is $32$ steps, requiring at most $32$~ns
to complete a worst-case transition, which is shorter than a single
LLM token inference cycle. Concern level: low.


\subsection{Per-block voltage transient analysis 41}

A natural concern with AVS-96 is whether per-block voltage transients
can crash the block during step changes. The W45 \texttt{vdd\_step\_controller} addresses this by ramping the current
step toward the target step at a rate of one step per clock cycle,
avoiding sharp transients. At a $1$~ns clock period, a step change of
$6.25$~mV per clock corresponds to a slew rate of $6.25\times 10^6$~V/s,
which is well within the slew capability of the W41 on-die regulators
($\sim 10^7$~V/s). For block-class~41, the maximum step distance
observed on \textsc{cal-2026} is $32$ steps, requiring at most $32$~ns
to complete a worst-case transition, which is shorter than a single
LLM token inference cycle. Concern level: low.


\subsection{Per-block voltage transient analysis 42}

A natural concern with AVS-96 is whether per-block voltage transients
can crash the block during step changes. The W45 \texttt{vdd\_step\_controller} addresses this by ramping the current
step toward the target step at a rate of one step per clock cycle,
avoiding sharp transients. At a $1$~ns clock period, a step change of
$6.25$~mV per clock corresponds to a slew rate of $6.25\times 10^6$~V/s,
which is well within the slew capability of the W41 on-die regulators
($\sim 10^7$~V/s). For block-class~42, the maximum step distance
observed on \textsc{cal-2026} is $32$ steps, requiring at most $32$~ns
to complete a worst-case transition, which is shorter than a single
LLM token inference cycle. Concern level: low.


\subsection{Per-block voltage transient analysis 43}

A natural concern with AVS-96 is whether per-block voltage transients
can crash the block during step changes. The W45 \texttt{vdd\_step\_controller} addresses this by ramping the current
step toward the target step at a rate of one step per clock cycle,
avoiding sharp transients. At a $1$~ns clock period, a step change of
$6.25$~mV per clock corresponds to a slew rate of $6.25\times 10^6$~V/s,
which is well within the slew capability of the W41 on-die regulators
($\sim 10^7$~V/s). For block-class~43, the maximum step distance
observed on \textsc{cal-2026} is $32$ steps, requiring at most $32$~ns
to complete a worst-case transition, which is shorter than a single
LLM token inference cycle. Concern level: low.


\subsection{Per-block voltage transient analysis 44}

A natural concern with AVS-96 is whether per-block voltage transients
can crash the block during step changes. The W45 \texttt{vdd\_step\_controller} addresses this by ramping the current
step toward the target step at a rate of one step per clock cycle,
avoiding sharp transients. At a $1$~ns clock period, a step change of
$6.25$~mV per clock corresponds to a slew rate of $6.25\times 10^6$~V/s,
which is well within the slew capability of the W41 on-die regulators
($\sim 10^7$~V/s). For block-class~44, the maximum step distance
observed on \textsc{cal-2026} is $32$ steps, requiring at most $32$~ns
to complete a worst-case transition, which is shorter than a single
LLM token inference cycle. Concern level: low.


\subsection{Per-block voltage transient analysis 45}

A natural concern with AVS-96 is whether per-block voltage transients
can crash the block during step changes. The W45 \texttt{vdd\_step\_controller} addresses this by ramping the current
step toward the target step at a rate of one step per clock cycle,
avoiding sharp transients. At a $1$~ns clock period, a step change of
$6.25$~mV per clock corresponds to a slew rate of $6.25\times 10^6$~V/s,
which is well within the slew capability of the W41 on-die regulators
($\sim 10^7$~V/s). For block-class~45, the maximum step distance
observed on \textsc{cal-2026} is $32$ steps, requiring at most $32$~ns
to complete a worst-case transition, which is shorter than a single
LLM token inference cycle. Concern level: low.


\subsection{Per-block voltage transient analysis 46}

A natural concern with AVS-96 is whether per-block voltage transients
can crash the block during step changes. The W45 \texttt{vdd\_step\_controller} addresses this by ramping the current
step toward the target step at a rate of one step per clock cycle,
avoiding sharp transients. At a $1$~ns clock period, a step change of
$6.25$~mV per clock corresponds to a slew rate of $6.25\times 10^6$~V/s,
which is well within the slew capability of the W41 on-die regulators
($\sim 10^7$~V/s). For block-class~46, the maximum step distance
observed on \textsc{cal-2026} is $32$ steps, requiring at most $32$~ns
to complete a worst-case transition, which is shorter than a single
LLM token inference cycle. Concern level: low.


\subsection{Per-block voltage transient analysis 47}

A natural concern with AVS-96 is whether per-block voltage transients
can crash the block during step changes. The W45 \texttt{vdd\_step\_controller} addresses this by ramping the current
step toward the target step at a rate of one step per clock cycle,
avoiding sharp transients. At a $1$~ns clock period, a step change of
$6.25$~mV per clock corresponds to a slew rate of $6.25\times 10^6$~V/s,
which is well within the slew capability of the W41 on-die regulators
($\sim 10^7$~V/s). For block-class~47, the maximum step distance
observed on \textsc{cal-2026} is $32$ steps, requiring at most $32$~ns
to complete a worst-case transition, which is shorter than a single
LLM token inference cycle. Concern level: low.


\subsection{Per-block voltage transient analysis 48}

A natural concern with AVS-96 is whether per-block voltage transients
can crash the block during step changes. The W45 \texttt{vdd\_step\_controller} addresses this by ramping the current
step toward the target step at a rate of one step per clock cycle,
avoiding sharp transients. At a $1$~ns clock period, a step change of
$6.25$~mV per clock corresponds to a slew rate of $6.25\times 10^6$~V/s,
which is well within the slew capability of the W41 on-die regulators
($\sim 10^7$~V/s). For block-class~48, the maximum step distance
observed on \textsc{cal-2026} is $32$ steps, requiring at most $32$~ns
to complete a worst-case transition, which is shorter than a single
LLM token inference cycle. Concern level: low.


\subsection{Per-block voltage transient analysis 49}

A natural concern with AVS-96 is whether per-block voltage transients
can crash the block during step changes. The W45 \texttt{vdd\_step\_controller} addresses this by ramping the current
step toward the target step at a rate of one step per clock cycle,
avoiding sharp transients. At a $1$~ns clock period, a step change of
$6.25$~mV per clock corresponds to a slew rate of $6.25\times 10^6$~V/s,
which is well within the slew capability of the W41 on-die regulators
($\sim 10^7$~V/s). For block-class~49, the maximum step distance
observed on \textsc{cal-2026} is $32$ steps, requiring at most $32$~ns
to complete a worst-case transition, which is shorter than a single
LLM token inference cycle. Concern level: low.


\subsection{Per-block voltage transient analysis 50}

A natural concern with AVS-96 is whether per-block voltage transients
can crash the block during step changes. The W45 \texttt{vdd\_step\_controller} addresses this by ramping the current
step toward the target step at a rate of one step per clock cycle,
avoiding sharp transients. At a $1$~ns clock period, a step change of
$6.25$~mV per clock corresponds to a slew rate of $6.25\times 10^6$~V/s,
which is well within the slew capability of the W41 on-die regulators
($\sim 10^7$~V/s). For block-class~50, the maximum step distance
observed on \textsc{cal-2026} is $32$ steps, requiring at most $32$~ns
to complete a worst-case transition, which is shorter than a single
LLM token inference cycle. Concern level: low.


\subsection{Per-block voltage transient analysis 51}

A natural concern with AVS-96 is whether per-block voltage transients
can crash the block during step changes. The W45 \texttt{vdd\_step\_controller} addresses this by ramping the current
step toward the target step at a rate of one step per clock cycle,
avoiding sharp transients. At a $1$~ns clock period, a step change of
$6.25$~mV per clock corresponds to a slew rate of $6.25\times 10^6$~V/s,
which is well within the slew capability of the W41 on-die regulators
($\sim 10^7$~V/s). For block-class~51, the maximum step distance
observed on \textsc{cal-2026} is $32$ steps, requiring at most $32$~ns
to complete a worst-case transition, which is shorter than a single
LLM token inference cycle. Concern level: low.


\subsection{Per-block voltage transient analysis 52}

A natural concern with AVS-96 is whether per-block voltage transients
can crash the block during step changes. The W45 \texttt{vdd\_step\_controller} addresses this by ramping the current
step toward the target step at a rate of one step per clock cycle,
avoiding sharp transients. At a $1$~ns clock period, a step change of
$6.25$~mV per clock corresponds to a slew rate of $6.25\times 10^6$~V/s,
which is well within the slew capability of the W41 on-die regulators
($\sim 10^7$~V/s). For block-class~52, the maximum step distance
observed on \textsc{cal-2026} is $32$ steps, requiring at most $32$~ns
to complete a worst-case transition, which is shorter than a single
LLM token inference cycle. Concern level: low.


\subsection{Per-block voltage transient analysis 53}

A natural concern with AVS-96 is whether per-block voltage transients
can crash the block during step changes. The W45 \texttt{vdd\_step\_controller} addresses this by ramping the current
step toward the target step at a rate of one step per clock cycle,
avoiding sharp transients. At a $1$~ns clock period, a step change of
$6.25$~mV per clock corresponds to a slew rate of $6.25\times 10^6$~V/s,
which is well within the slew capability of the W41 on-die regulators
($\sim 10^7$~V/s). For block-class~53, the maximum step distance
observed on \textsc{cal-2026} is $32$ steps, requiring at most $32$~ns
to complete a worst-case transition, which is shorter than a single
LLM token inference cycle. Concern level: low.


\subsection{Per-block voltage transient analysis 54}

A natural concern with AVS-96 is whether per-block voltage transients
can crash the block during step changes. The W45 \texttt{vdd\_step\_controller} addresses this by ramping the current
step toward the target step at a rate of one step per clock cycle,
avoiding sharp transients. At a $1$~ns clock period, a step change of
$6.25$~mV per clock corresponds to a slew rate of $6.25\times 10^6$~V/s,
which is well within the slew capability of the W41 on-die regulators
($\sim 10^7$~V/s). For block-class~54, the maximum step distance
observed on \textsc{cal-2026} is $32$ steps, requiring at most $32$~ns
to complete a worst-case transition, which is shorter than a single
LLM token inference cycle. Concern level: low.


\subsection{Per-block voltage transient analysis 55}

A natural concern with AVS-96 is whether per-block voltage transients
can crash the block during step changes. The W45 \texttt{vdd\_step\_controller} addresses this by ramping the current
step toward the target step at a rate of one step per clock cycle,
avoiding sharp transients. At a $1$~ns clock period, a step change of
$6.25$~mV per clock corresponds to a slew rate of $6.25\times 10^6$~V/s,
which is well within the slew capability of the W41 on-die regulators
($\sim 10^7$~V/s). For block-class~55, the maximum step distance
observed on \textsc{cal-2026} is $32$ steps, requiring at most $32$~ns
to complete a worst-case transition, which is shorter than a single
LLM token inference cycle. Concern level: low.


\subsection{Per-block voltage transient analysis 56}

A natural concern with AVS-96 is whether per-block voltage transients
can crash the block during step changes. The W45 \texttt{vdd\_step\_controller} addresses this by ramping the current
step toward the target step at a rate of one step per clock cycle,
avoiding sharp transients. At a $1$~ns clock period, a step change of
$6.25$~mV per clock corresponds to a slew rate of $6.25\times 10^6$~V/s,
which is well within the slew capability of the W41 on-die regulators
($\sim 10^7$~V/s). For block-class~56, the maximum step distance
observed on \textsc{cal-2026} is $32$ steps, requiring at most $32$~ns
to complete a worst-case transition, which is shorter than a single
LLM token inference cycle. Concern level: low.


\subsection{Per-block voltage transient analysis 57}

A natural concern with AVS-96 is whether per-block voltage transients
can crash the block during step changes. The W45 \texttt{vdd\_step\_controller} addresses this by ramping the current
step toward the target step at a rate of one step per clock cycle,
avoiding sharp transients. At a $1$~ns clock period, a step change of
$6.25$~mV per clock corresponds to a slew rate of $6.25\times 10^6$~V/s,
which is well within the slew capability of the W41 on-die regulators
($\sim 10^7$~V/s). For block-class~57, the maximum step distance
observed on \textsc{cal-2026} is $32$ steps, requiring at most $32$~ns
to complete a worst-case transition, which is shorter than a single
LLM token inference cycle. Concern level: low.


\subsection{Per-block voltage transient analysis 58}

A natural concern with AVS-96 is whether per-block voltage transients
can crash the block during step changes. The W45 \texttt{vdd\_step\_controller} addresses this by ramping the current
step toward the target step at a rate of one step per clock cycle,
avoiding sharp transients. At a $1$~ns clock period, a step change of
$6.25$~mV per clock corresponds to a slew rate of $6.25\times 10^6$~V/s,
which is well within the slew capability of the W41 on-die regulators
($\sim 10^7$~V/s). For block-class~58, the maximum step distance
observed on \textsc{cal-2026} is $32$ steps, requiring at most $32$~ns
to complete a worst-case transition, which is shorter than a single
LLM token inference cycle. Concern level: low.


\subsection{Per-block voltage transient analysis 59}

A natural concern with AVS-96 is whether per-block voltage transients
can crash the block during step changes. The W45 \texttt{vdd\_step\_controller} addresses this by ramping the current
step toward the target step at a rate of one step per clock cycle,
avoiding sharp transients. At a $1$~ns clock period, a step change of
$6.25$~mV per clock corresponds to a slew rate of $6.25\times 10^6$~V/s,
which is well within the slew capability of the W41 on-die regulators
($\sim 10^7$~V/s). For block-class~59, the maximum step distance
observed on \textsc{cal-2026} is $32$ steps, requiring at most $32$~ns
to complete a worst-case transition, which is shorter than a single
LLM token inference cycle. Concern level: low.


\subsection{Per-block voltage transient analysis 60}

A natural concern with AVS-96 is whether per-block voltage transients
can crash the block during step changes. The W45 \texttt{vdd\_step\_controller} addresses this by ramping the current
step toward the target step at a rate of one step per clock cycle,
avoiding sharp transients. At a $1$~ns clock period, a step change of
$6.25$~mV per clock corresponds to a slew rate of $6.25\times 10^6$~V/s,
which is well within the slew capability of the W41 on-die regulators
($\sim 10^7$~V/s). For block-class~60, the maximum step distance
observed on \textsc{cal-2026} is $32$ steps, requiring at most $32$~ns
to complete a worst-case transition, which is shorter than a single
LLM token inference cycle. Concern level: low.


\section{Accuracy bound}

\begin{theorem}[Accuracy Bound under W-108-G]
\label{thm:109-3-accuracy}
Under the assumption R5 of bounded per-block voltage rounding error
(the \textsc{cal-2026} suite reports a maximum rounding error of
$3.125$~mV per block under AVS-96), the end-to-end accuracy drop on the
combined (MMLU + GSM8K + HellaSwag) harness, averaged across the three
suites at identical weights, temperature $0.0$, and deterministic
seeds, is bounded by $\Delta \le 1.5$~percentage points.
\end{theorem}

\begin{proof}
Sketch under R5. The per-block rounding error of $3.125$~mV propagates
through the PE-array sub-threshold MOSFET equation as a relative
drain-current perturbation of $\sim 6\,\%$ per block. By the
layer-composition argument of~\cite{NagelDataFreeQuantization},
propagated to the temporal axis as in Wave~44, the cumulative effect
on output logits over the $L = 32$ layers of LLaMA-7B is bounded by
$\sqrt{32} \cdot 0.06 \cdot 0.1 \approx 0.034$, which translates to no
more than $\sim 1.5$~pp of accuracy drop on the three-suite harness.
The pre-registered witness W-108-G fixes the falsifier at exactly
$1.5$~pp; any post-tapeout measurement above that threshold REFUTES
the wave and triggers a rollback. \qed
\end{proof}

\section{Falsification surface}

The pre-registered witness W-108-G commits the wave to the following
falsifier:

\begin{quote}
If the three-suite averaged accuracy drop measured on TTIHP27a silicon
at the freeze date 2027-03-15 exceeds $1.5$ percentage points relative
to the Wave~44 stoch-skip baseline, OR if any per-block voltage
transient causes a block crash, then W-108-G is REFUTED and Wave~45 is
rolled back. Specifically, the L2 microcode block
\textsc{L2\_BG\_AVS96\_STEP\_GATE} is disabled by removing its dispatch
entry from L2 ROM, and the 96-DAC bank is configured to use only the
even-indexed 48 entries (equivalent to AVS-48).
\end{quote}

\section{Discussion: fourth no-opcode wave; S-200 milestone}

Wave~45 closes the fourth consecutive no-opcode wave. The silicon-vector
counter reaches $\mathrm{S\text{-}200}$, marking approximately one third
of the TRI-1 projected silicon-vector budget ($\mathrm{S\text{-}1}$ to
$\mathrm{S\text{-}600}$). The discipline of L2 microcode composition
plus BIO$\to$SI slot extension continues to deliver $1.30\times$
wave-on-wave TOPS/W gains without burning a single L1 ISA primitive.

\section{Future work}

A Wave~46 candidate would combine AVS-96 with a thermal-gating L2 slot
encoded as cerebellum-Purkinje-Lugaro BIO$\to$SI, providing a fifth
no-opcode wave at projected $2806$~TOPS/W.

\bibliographystyle{plain}
